\documentclass[12pt]{article}
\usepackage[utf8]{inputenc}
\usepackage[T1]{fontenc}
\usepackage[margin=1in]{geometry}
\usepackage{setspace}
\usepackage{xcolor}
\usepackage[colorlinks=true, urlcolor=blue!60!black, linkcolor=blue!60!black]{hyperref}
\setlength{\parindent}{2em}
\setlength{\parskip}{1.25em}
\renewcommand{\baselinestretch}{1.0}
\title{\textit{A.L.A. Schechter Poultry Corp. v. United States}, 295 U.S. 495 (1935)}
\author{Case Brief}
\date{}

\begin{document}
\maketitle

\section{Case Metadata}
\begin{description}
  \item[Citation:] 295 U.S. 495 (1935) (Nos.\ 854 \& 864; consolidated with \textit{United States v. A.L.A. Schechter Poultry Corp.}).
  \item[Argued / Decided:] Argued May 2--3, 1935; Decided May 27, 1935.
  \item[Vote:] 9--0 (Unanimous).
  \item[Author:] Chief Justice Charles Evans Hughes (joined by Justices Van Devanter, McReynolds, Brandeis, Sutherland, Butler, and Roberts; Justices Cardozo and Stone joined the judgment and concurred separately).
  \item[Concurrences / Dissents:] Justice Benjamin N.\ Cardozo filed a concurring opinion, in which Justice Harlan Fiske Stone joined. No dissents.
  \item[Chapter / Topic:] Separation of Powers and the Executive (Nondelegation Doctrine; Article I Legislative Power; Commerce Clause).
  \item[Primary Source URLs:] \href{https://supreme.justia.com/cases/federal/us/295/495/}{Justia U.S. Supreme Court Center} (\url{https://supreme.justia.com/cases/federal/us/295/495/}); \href{https://www.law.cornell.edu/supremecourt/text/295/495}{Cornell LII} (\url{https://www.law.cornell.edu/supremecourt/text/295/495}).
\end{description}

\section{Facts and Procedural History}
At the height of the Great Depression, Congress enacted the National Industrial Recovery Act (NIRA) of June 16, 1933 (48 Stat.\ 195). Section~3 of Title~I authorized the President, upon application by ``truly representative'' trade or industrial associations, to approve ``codes of fair competition''---or to prescribe them on his own motion---so long as the codes did not promote monopolies or oppress small enterprises and would ``tend to effectuate the policy of this title'' as declared in \S\,1. Under \S\,3(f), violating any provision of an approved code ``in any transaction in or affecting interstate or foreign commerce'' was a federal misdemeanor punishable by a fine of up to \$500 per day.

On April 13, 1934, President Franklin D.\ Roosevelt approved the ``Code of Fair Competition for the Live Poultry Industry of the Metropolitan Area in and about the City of New York,'' governing live poultry commerce across New York City and surrounding counties in New York, New Jersey, and Connecticut. The Code established a 40-hour workweek, a 50-cent hourly minimum wage, child-labor prohibitions, collective-bargaining rights, mandatory price and sales reporting, and trade-practice rules. Chief among the trade-practice rules was the ``straight killing'' requirement (Art.\ VII, \S\,14), which barred wholesale slaughterhouses from allowing retail buyers to select individual chickens from coops or half-coops, requiring buyers instead to accept the run of the coop (except for culls).

A.L.A. Schechter Poultry Corp.\ and Schechter Live Poultry Market---wholesale kosher poultry slaughterhouses in Brooklyn operated by Joseph, Martin, Alex, and Aaron Schechter---bought live poultry from commission men at the West Washington Market in Manhattan or railroad terminals (and occasionally Philadelphia). Although 96\% of live poultry sold in New York City arrived from other states, the Schechters trucked their purchased birds to Brooklyn, slaughtered them via \textit{schochtim}, and sold them within 24 hours exclusively to local retail dealers and butchers; they sold no poultry in interstate commerce.

A federal grand jury in the Eastern District of New York indicted the Schechters on 60 counts. The trial court dismissed 27 counts, a jury acquitted on 14, and the defendants were convicted on 19 counts: one conspiracy count and 18 substantive code violations, including two wage-and-hour counts, ten ``straight killing'' counts, one count for selling an unfit chicken (giving rise to the case's nickname, the ``sick chicken case''), two counts for violating municipal inspection ordinances, two counts for false or omitted price/volume reports, and one count for selling to unlicensed dealers. The Second Circuit (76 F.2d 617) reversed the two wage-and-hour convictions as exceeding Congress's Commerce Clause power, but affirmed the conspiracy count and the remaining 16 substantive counts. Both the Schechters (No.\ 854) and the United States (No.\ 864) sought review, and the Supreme Court granted cross-writs of certiorari.

\section{Issues Presented}
\begin{enumerate}
  \item \textbf{Nondelegation (Article I, \S\,1):} Does \S\,3 of the National Industrial Recovery Act unconstitutionally delegate legislative power to the President and private trade associations by authorizing the promulgation of binding ``codes of fair competition'' without intelligible statutory standards?
  \item \textbf{Commerce Clause (Article I, \S\,8, cl.\ 3):} Do the Live Poultry Code's wage, hour, and local-sales provisions exceed Congress's power to regulate interstate commerce as applied to intrastate poultry slaughtering and local retail sales?
\end{enumerate}

\section{Holding}
Yes on both issues. Section~3 of the NIRA is an unconstitutional delegation of legislative power under Article~I, \S\,1, and the Live Poultry Code's application to the Schechters' intrastate slaughtering and sales exceeds Congress's power under the Commerce Clause. The Court reversed the remaining convictions (No.\ 854) and affirmed the Second Circuit's reversal of the wage-and-hour counts (No.\ 864), declining to reach the Fifth Amendment Due Process claim.

\section{Reasoning}
Writing for a unanimous Court, Chief Justice Hughes first rejected the Government's threshold arguments. While the Great Depression was a grave national crisis, ``[e]xtraordinary conditions do not create or enlarge constitutional power'' (295 U.S.\ at 528). Nor were the NIRA codes merely voluntary cooperative arrangements: once approved, they functioned as ``codes of laws'' backed by criminal sanctions, binding assenting and non-assenting businesses alike.

\begin{enumerate}
  \item \textbf{Unconstitutional Delegation of Legislative Power:} Grounding its analysis in Article~I, \S\,1 and the Necessary and Proper Clause (Art.\ I, \S\,8, cl.\ 18), and building on \textit{Panama Refining Co. v. Ryan}, 293 U.S.\ 388 (1935), the Court reiterated that Congress may not abdicate or transfer its essential legislative functions. Although Congress may enlist executive instrumentalities to make subordinate rules within prescribed limits and find facts to trigger declared legislative policy, Congress itself must lay down the policies and establish standards.

  Unlike \textit{Panama Refining}---where Congress at least defined the prohibited subject matter (``hot oil'')---NIRA \S\,3 failed to define ``fair competition.'' The Court demonstrated that ``fair competition'' was neither the common-law antithesis of ``unfair competition'' (palming off or misappropriation) nor analogous to ``unfair methods of competition'' under \S\,5 of the Federal Trade Commission Act, which was cabined by quasi-judicial agency procedures, notice, hearings, evidentiary findings, and judicial review. Instead, \S\,3 swept in whatever rules trade associations proposed and the President approved to achieve the sweeping, conflicting aspirations in \S\,1's Declaration of Policy (rehabilitating industry, relieving unemployment, increasing purchasing power, conserving resources, and promoting cooperative action).

  The Court rejected any suggestion that Congress could delegate lawmaking power to private trade associations as ``unknown to our law'' and ``utterly inconsistent with the constitutional prerogatives and duties of Congress'' (295 U.S.\ at 537). Turning to the President's approval authority, Hughes found the statutory constraints illusory: the requirement that the President find a code ``will tend to effectuate the policy of this title'' was merely a ``statement of an opinion,'' and the President could add, modify, or prescribe codes on his own motion without being bound by any administrative agency's findings. Distinguishing upheld delegations to the Interstate Commerce Commission, the Federal Radio Commission (\textit{FRC v. Nelson Bros.\ Co.}, 289 U.S.\ 266 (1933)), and the President's flexible tariff authority (\textit{J.W. Hampton, Jr., \& Co. v. United States}, 276 U.S.\ 394 (1928)), Hughes concluded that \S\,3 was ``without precedent,'' leaving the President ``virtually unfettered'' discretion to enact laws governing the entire national economy.

  \item \textbf{Commerce Clause Limits (Direct vs.\ Indirect Effects):} Independently, the Court held that the Live Poultry Code exceeded federal power under the Commerce Clause. First, the poultry was no longer ``in'' interstate commerce or in a ``stream'' or ``current'' of commerce (distinguishing \textit{Swift \& Co. v. United States}, 196 U.S.\ 375 (1905), and \textit{Stafford v. Wallace}, 258 U.S.\ 495 (1922)). Once the birds reached the Schechters' Brooklyn slaughterhouses for local slaughter and sale, interstate movement had ended and the goods had come to ``a permanent rest within the State'' (295 U.S.\ at 543).

  Second, the Schechters' local wages, hours, and ``straight killing'' practices did not \textit{directly} affect interstate commerce. Drawing a sharp doctrinal line between ``direct'' and ``indirect'' effects---a distinction Hughes called ``fundamental'' and ``essential to the maintenance of our constitutional system''---the Court warned that if Congress could regulate local labor costs and sales merely because they affected prices and indirectly influenced interstate markets, federal power would know no bounds and the States' authority over domestic concerns would exist ``only by sufferance of the federal government'' (295 U.S.\ at 546, 548).
\end{enumerate}

\section{Concurring and Dissenting Opinions}
\textbf{Justice Cardozo, concurring (joined by Justice Stone):} Having been the lone dissenter in \textit{Panama Refining} just four months earlier, Cardozo wrote separately to explain why NIRA \S\,3 crossed the constitutional line where \S\,9(c) had not. In \textit{Panama Refining}, Congress had conferred discretion only over the \textit{occasion} for prohibiting a single, defined act (transporting excess oil), not a ``roving commission.'' Here, by contrast, the power was ``not canalized within banks that keep it from overflowing''; it was ``unconfined and vagrant''---a ``roving commission to inquire into evils and, upon discovery, correct them'' (295 U.S.\ at 551). Because ``codes of fair competition'' encompassed positive economic planning rather than merely suppressing recognized unfair trade practices, Cardozo famously declared: ``This is delegation running riot'' (295 U.S.\ at 553). On the Commerce Clause, Cardozo acknowledged Judge Learned Hand's observation below that society is an ``elastic medium which transmits all tremors throughout its territory,'' but insisted that ``[t]he law is not indifferent to considerations of degree''; local activities do not become interstate merely through distant repercussions, lest ``there will be an end to our federal system'' (295 U.S.\ at 554).

\section{Key Quotations}
\begin{quote}
``Extraordinary conditions may call for extraordinary remedies. But the argument necessarily stops short of an attempt to justify action which lies outside the sphere of constitutional authority. Extraordinary conditions do not create or enlarge constitutional power.''\\
--- Chief Justice Hughes, 295 U.S.\ at 528.
\end{quote}

\begin{quote}
``The Congress is not permitted to abdicate or to transfer to others the essential legislative functions with which it is thus vested.''\\
--- Chief Justice Hughes, 295 U.S.\ at 529.
\end{quote}

\begin{quote}
``Congress cannot delegate legislative power to the President to exercise an unfettered discretion to make whatever laws he thinks may be needed or advisable for the rehabilitation and expansion of trade or industry.''\\
--- Chief Justice Hughes, 295 U.S.\ at 537--538.
\end{quote}

\begin{quote}
``Section 3 of the Recovery Act is without precedent. It supplies no standards for any trade, industry or activity. It does not undertake to prescribe rules of conduct to be applied to particular states of fact determined by appropriate administrative procedure. Instead of prescribing rules of conduct, it authorizes the making of codes to prescribe them.''\\
--- Chief Justice Hughes, 295 U.S.\ at 541.
\end{quote}

\begin{quote}
``The delegated power of legislation which has found expression in this code is not canalized within banks that keep it from overflowing. It is unconfined and vagrant \dots\ This is delegation running riot. No such plenitude of power is susceptible of transfer.''\\
--- Justice Cardozo, concurring, 295 U.S.\ at 551, 553.
\end{quote}

\section{Doctrinal Significance and Modern Legacy (Through 2026)}
Decided on ``Black Monday'' (May 27, 1935) alongside \textit{Humphrey's Executor v. United States}, 295 U.S.\ 602 (1935), \textit{Schechter Poultry} dismantled the centerpiece of the early New Deal and remains one of the two most famous structural decisions in American constitutional law---yet its two holdings experienced sharply divergent fates:

\begin{itemize}
  \item \textbf{Commerce Clause Repudiation and Partial Revival:} Within seven years, the Court abandoned \textit{Schechter}'s formalist direct/indirect effects test in \textit{NLRB v. Jones \& Laughlin Steel Corp.}, 301 U.S.\ 1 (1937), \textit{United States v. Darby}, 312 U.S.\ 100 (1941), and \textit{Wickard v. Filburn}, 317 U.S.\ 111 (1942), adopting the modern ``substantial effects'' and aggregation framework. Nevertheless, Hughes's warning that the Commerce Clause must not erase the boundary between national and local power resurfaced in \textit{United States v. Lopez}, 514 U.S.\ 549 (1995), and \textit{NFIB v. Sebelius}, 567 U.S.\ 519 (2012).
  \item \textbf{The ``One Good Year'' of Facial Nondelegation:} Together with \textit{Panama Refining}, \textit{Schechter Poultry} remains the \textit{only} Supreme Court decision ever to strike down a federal statute under the Article~I nondelegation doctrine. Throughout the post-1935 era, the Court applied the lenient ``intelligible principle'' test of \textit{J.W. Hampton} to uphold sweeping delegations in \textit{Yakus v. United States}, 321 U.S.\ 414 (1944), \textit{Mistretta v. United States}, 488 U.S.\ 361 (1989), and \textit{Whitman v. American Trucking Ass'ns}, 531 U.S.\ 457 (2001).
  \item \textbf{\textit{Gundy v. United States}, 588 U.S.\ 128 (2019)} (\href{https://supreme.justia.com/cases/federal/us/588/17-6086/}{Justia}): A four-Justice plurality (Kagan, J.) narrowly construed SORNA \S\,20913(d) to survive review, citing \textit{Schechter} and \textit{Panama Refining} as the lone historical outliers. Justice Gorsuch's dissent (joined by Roberts, C.J., and Thomas, J.) urged replacing the modern intelligible-principle test with an originalist framework limiting executive discretion to filling up details, fact-finding, and non-legislative functions, while Justice Alito concurred in the judgment to signal openness to reconsidering the doctrine with a full Court.
  \item \textbf{Sub-Constitutional Enforcement via Major Questions and \textit{Loper Bright}:} Rather than reviving facial invalidations under \textit{Schechter}, the Roberts Court channeled nondelegation values into statutory clear-statement rules---specifically the \textbf{major questions doctrine} (\textit{West Virginia v. EPA}, 597 U.S.\ 697 (2022); \textit{Biden v. Nebraska}, 600 U.S.\ 477 (2023))---and overruled \textit{Chevron} deference in \textit{Loper Bright Enterprises v. Raimondo}, 603 U.S.\ 369 (2024).
  \item \textbf{\textit{FCC v. Consumers' Research}, 606 U.S.\ 656 (2025)} (\href{https://supreme.justia.com/cases/federal/us/606/24-354/}{Justia}; \href{https://www.supremecourt.gov/opinions/24pdf/24-354_0861.pdf}{Slip Op.}): In a 6--3 decision (Kagan, J., joined by Roberts, C.J., and Sotomayor, Kavanaugh, Barrett, and Jackson, JJ.), the Court reversed the en banc Fifth Circuit and upheld the FCC's Universal Service Fund contribution scheme under \S\,254 of the 1996 Telecommunications Act. The majority reaffirmed the intelligible-principle test, rejected a stricter numeric-cap rule for revenue-raising delegations, upheld the advisory role of the private Universal Service Administrative Company under \textit{Sunshine Anthracite Coal Co. v. Adkins}, 310 U.S.\ 381 (1940) (distinguishing \textit{Schechter} and \textit{Carter Coal}), and rejected the Fifth Circuit's ``combination'' theory. Justice Kagan expressly invoked \textit{Schechter} (295 U.S.\ at 521--522, 541--542) as an ``object lesson'' of virtually unfettered discretion, while Justice Kavanaugh's concurrence emphasized that \textit{Loper Bright} and the major questions doctrine now do much of the structural work---though warning that delegations to independent agencies insulated from presidential removal raise unresolved Article~II concerns. Justice Gorsuch's dissent (joined by Thomas and Alito, JJ.) cited \textit{Schechter} (295 U.S.\ at 534, 537, 539--542) to argue that \S\,254's competing goals were a ``preface of generalities'' unconstitutionally delegating the taxing power.
  \item \textbf{\textit{Learning Resources, Inc. v. Trump}, 607 U.S.\ 229 (2026)} (\href{https://supreme.justia.com/cases/federal/us/607/24-1287/}{Justia}; \href{https://www.supremecourt.gov/opinions/25pdf/24-1287_4gcj.pdf}{Slip Op.}): In a 6--3 ruling (Roberts, C.J., writing for the Court in Parts~I, II-A-1, and II-B, and for a three-Justice plurality with Gorsuch and Barrett, JJ., in Parts~II-A-2 and III; Kagan, J., joined by Sotomayor and Jackson, JJ., concurring in part and in the judgment), the Court held that the International Emergency Economic Powers Act (IEEPA) does not authorize the President to impose tariffs. Echoing \textit{Schechter}'s warning against using emergencies to justify unbounded executive lawmaking, the plurality applied the major questions doctrine to protect Congress's core Article~I taxing power, demonstrating that in 2026 \textit{Schechter}'s anti-abdication principle lives primarily through clear-statement statutory review rather than facial invalidations.
\end{itemize}

\section{Pedagogical Critiques and Exam Hypotheticals}
\textbf{Pedagogical Critiques:}
\begin{itemize}
  \item \textit{The Private-Delegation Dimension:} Scholars note that \textit{Schechter} was not solely about executive discretion; it also reflected deep suspicion of corporatist self-regulation---handing coercive lawmaking power to private trade associations whose dominant members could cartelize industries and crush smaller competitors like the Schechter brothers (a theme cemented a year later in \textit{Carter v. Carter Coal Co.}, 298 U.S.\ 238 (1936), and revisited in \textit{FCC v. Consumers' Research} (2025)).
  \item \textit{Judicial Manageability:} Justice Scalia famously observed in \textit{Mistretta} (488 U.S.\ at 415--416, dissenting) that while the nondelegation principle is constitutionally sound, determining ``how much is too much'' is a question of degree poorly suited to bright-line judicial enforcement---explaining why courts prefer statutory narrowing (\textit{Gundy}, \textit{Learning Resources}) over \textit{Schechter}-style invalidation.
  \item \textit{Procedural Safeguards as a Proxy for Standards:} Notice how Chief Justice Hughes contrasted NIRA \S\,3 with the FTC Act and ICC statutes: the absence of notice-and-comment rulemaking, formal hearings, evidentiary records, and judicial review (later standardized by the 1946 Administrative Procedure Act) made NIRA's delegation uniquely dangerous.
\end{itemize}

\textbf{``Opposite Side'' Contrast (\textit{Schechter Poultry} vs.\ \textit{Curtiss-Wright} and \textit{Consumers' Research}):}
Contrast \textit{Schechter Poultry} (1935) with \textit{United States v. Curtiss-Wright Export Corp.}, 299 U.S.\ 304 (1936), decided just 19 months later by nearly the exact same Justices (with only Stone recused and McReynolds dissenting). In \textit{Schechter}, domestic economic regulation across all U.S.\ industries demanded clear congressional standards and procedural guardrails; in \textit{Curtiss-Wright}, an arms embargo aimed at a foreign war received sweeping judicial deference because of the President's independent Article~II foreign-affairs role. Similarly, contrast \textit{Schechter} with \textit{FCC v. Consumers' Research}, 606 U.S.\ 656 (2025): whereas NIRA \S\,3 allowed private industry groups and the President to invent substantive criminal codes across every sector of the economy under conflicting generalities, \S\,254 of the 1996 Telecommunications Act targeted one sector, specified statutory beneficiaries and mandatory criteria, and kept the private administrator strictly subordinate to agency oversight.

\textbf{Exam Hypothetical:}
Suppose Congress enacts the ``Artificial Intelligence Stabilization Act of 2026,'' authorizing the President to approve ``Codes of Safe AI Competition'' drafted by a consortium of major AI laboratories whenever the President finds a code ``tends to promote innovation, national competitiveness, and public safety.'' Violations carry criminal fines up to \$100,000 per day, and the President's approval is exempted from the Administrative Procedure Act and judicial review. A small open-source AI startup indicted for violating a code's compute-reporting rule challenges the statute.
\begin{itemize}
  \item \textit{Challenger's Best Arguments:} Directly tracks \textit{Schechter Poultry} and \textit{Carter Coal}: (1) standardless, conflicting statutory aims (``innovation'' vs.\ ``public safety''); (2) private delegation to self-interested industry incumbents with no binding agency oversight (distinguishing \textit{Consumers' Research} and \textit{Sunshine Anthracite}); (3) elimination of APA procedural safeguards and judicial review; and (4) criminal penalties across a transformative economic sector (triggering both \textit{Schechter} and, if ambiguous, the major questions doctrine under \textit{West Virginia v. EPA} and \textit{Learning Resources}).
  \item \textit{Government's Best Arguments:} Argue that ``public safety'' and ``innovation'' supply an intelligible principle comparable to ``public interest, convenience, and necessity'' (\textit{NBC}, \textit{Consumers' Research}); argue that the AI consortium is strictly advisory because only presidential approval gives the code legal force (citing Cardozo's concurrence in \textit{Schechter} and \textit{Consumers' Research}); and invoke national-security/foreign-competition aspects of AI to seek \textit{Curtiss-Wright} deference (though challengers will reply under \textit{Learning Resources} (2026) that domestic regulatory codes fall within Congress's exclusive Article~I sphere).
\end{itemize}

\end{document}
